\documentclass{article}
\usepackage[a4paper, total={6.5in, 10in}]{geometry}
\setlength{\parindent}{0pt}
\usepackage{tcolorbox}
\tcbset{colback=white!94!black, colframe=black!60!white, coltitle = white, boxrule=0.8pt, arc=2mm}
\newtcolorbox{boxs}[2][]{title= #2,#1}
\usepackage{datetime2}
\usepackage[british]{babel}
\usepackage{amsfonts}
\usepackage{amsmath}

\title{\textbf{Advanced Calculus HW 8}}
\author{Robert O'Connell}
\date{\today}
\begin{document}
\maketitle
\vspace{5mm}
\subsection*{Question 1}
\subsubsection*{(a)}
\[
f(x,y) = e^{-x^2 - y^2 + 4x}
\]
\[
\nabla f(x.y) = ((-2+4)e^{-x^2 - y^2 + 4x}, (-2y)e^{-x^2 - y^2 + 4x})
\]
\[
\nabla f(x,y) = \vec{0} \iff (x,y) = (2,0)
\]
Hence \((2,0)\) is a critical point

Then using the second partials test, we need to compute
\[
D = f_{xx}(2,0) f_{yy}(2,0) - (f_{xy}(2,0))^2
\]

\[
f_{xx}(x,y) = -2e^{-x^2 - y^2 + 4x} + (-2x + 4)^2e^{-x^2 - y^2 + 4x}
\]
\[
f_{xx}(2,0) = -2e^4
\]

\[
f_{yy}(x,y) = -2e^{-x^2 - y^2 + 4x} + 4y^2e^{-x^2 - y^2 + 4x}
\]
\[
f_{yy}(2,0) = -2e^4
\]

\[
f_{xy}(x,y) = (-2x+4)(-2y)e^{-x^2 - y^2 + 4x}
\]
\[
f_{xy}(2,0) = 0
\]

Hence \(D = 4e^4 >0\)

Since \(f_{xx}(2,0)>0\) as well, we have that there is a relative maximum of \(f\) at \((2,0)\) 
\subsubsection*{(b)}
\[
f(x,y) = xy + \frac{2}{x} + \frac{4}{y}
\]
Note: \(f\) is not defined for all points with \(x=0\) or \(y=0\)
\[
\nabla f(x,y) = (y- \frac{2}{x^2}, x - \frac{4}{y^2})
\]
\[
\nabla f(x,y) = \vec{0} \iff x^2y =2 \text{ and } xy^2 = 4
\]
The only such point is \((1,2)\)

Hence \((1,2)\) is a critical point

Then using the second partials test, we need to compute
\[
D = f_{xx}(1,2) f_{yy}(1,2) - (f_{xy}(1,2))^2
\]

\[
f_{xx}(x,y) = \frac{4}{x^3}
\]
\[
f_{xx}(1,2) = 4
\]

\[
f_{yy}(x,y) = \frac{8}{y^3}
\]
\[
f_{yy}(1,2) = 1
\]

\[
f_{xy}(x,y) = 1
\]

Hence \(D = 3>0 \) 

Since \(f_{xx}(1,2)>0\) as well, we have that there is a relative maximum of \(f\) at \((1,2)\) 



\subsubsection*{(c)}
\[
f(x,y) = e^x \cos(y)
\]

\[
\nabla f(x,y) = (e^x \cos(y) , -e^x \sin(y))
\]
Since \(e^x \neq 0 \ \forall \ x\) and \(\cos(y) = \sin(y) \neq 0 \ \forall \ y\),

\(\nabla f(x,y) \neq \vec{0} \ \forall \ x,y\)

Hence, no critical points and hence no relative maxima, minima or saddle points

\subsection*{Question 2}
\subsubsection*{Step 1}
Find all interior critical points
\[
\nabla f(x,y) = (y^2, 2yx)
\]
\(\nabla f(x,y) = \vec{0} \iff y =0\) or \(x=0\), all such points are on the boundary, hence no interior critical points

\subsubsection*{Step 2}
Find all boundary points on which relative extrema can occur

\begin{itemize}
    \item \(x \in [0,1], y =0\) \\
    \(f(x,0) = 0\) \\
    All such points are candidates \\

    \item \(y \in [0,1], x=0\) \\
    \(f(0,y) = 0\) \\
    All such points are candidates \\

    \item All points s.t. \(x > 0, y> 0, x^2 + y^2 = 1\) \\
    \(x \in (0,1), y = \sqrt{1 - x^2}\) \\
    \(f(x,\sqrt{1-x^2}) = x(1-x^2) = x - x^3\) \\
    Extremised at \(f \ ' (x,\sqrt{1-x^2}) = 0 \iff 1 - 3x^2 = 0\) \\
    Only at \(x = \frac{1}{\sqrt{3}}\) \\
    Then at the point \(\left(\frac{1}{\sqrt{3}}, \sqrt{\frac{2}{3}}\right)\)
\end{itemize}

\subsubsection*{Step 3}
Evaluate all found points
\[
f(x,0) = 0 \ \forall \ x
\]
\[
f(0,y) =0 \ \forall \ y
\]
\[
f\left(\frac{1}{\sqrt{3}}, \sqrt{\frac{2}{3}}\right) = \frac{2}{3\sqrt{3}}
\]

Hence the absolute minima of this function on \(\mathcal{R}\) is 0 and the absolute maximum is \( \frac{2}{3\sqrt{3}}\)
\subsection*{Question 3}
As \(P(1,-1,1)\) lies on the plane \(4x + 3y +z =2\), the closest point on the plane to \(P\) is simply \(P\)

Using the method of Lagrange multipliers, we have the function to minimize as the square of a point to the point \(P(1,-1,1)\), 
\(f(x) = (x-1)^2 + (y+1)^2 + (z-1)^2\), and with \(g(x) = 4x + 3y +z -2\), the equation of the plane

Then we calculate 
\[
\nabla f(x,y,z) = (2(x-1), 2(y+1), 2(z-1))
\]
\[
\nabla g(x,y,z) = (4,3,1)
\]

Then, we find \(\lambda\) such that \(\nabla f = \lambda \nabla g\)
\[
(2(x-1), 2(y+1), 2(z-1)) = (4\lambda, 3\lambda, \lambda)
\]

Solving these three linear equations we get \((x,y,z) = (1,-1,1)\), as expected
\end{document}